% =============================================================================
% ISI Paper 3 -- Analytisches Briefing (German)
% Externe Lieferantenkonzentration in der EU-27
% Sections 1--7
% =============================================================================

% ─────────────────────────────────────────────────────────────────────────────
\section{Zusammenfassung}
\label{sec:summary}
% ─────────────────────────────────────────────────────────────────────────────

\begin{sloppypar}
\textbf{Kernaussage.}
Der International Sovereignty Index (\ISI) ist ein
sechs\-dimensionaler Kompositindikator der externen
Lieferanten\-konzentration.  Die empirische Anwendung auf die EU-27
(Datenjahrgang~2024) zeigt, dass die Querschnitts\-varianz der
Kompositwerte \"uberwiegend durch zwei der sechs Achsen bestimmt
wird: R\"ustung und Logistik.  Dieses Ergebnis ist ein deskriptiver
Varianzbefund --- keine konzeptionelle Reduktion des Index auf
zwei Dimensionen.
\end{sloppypar}

\medskip
\textbf{Empirische Kernergebnisse:}
\begin{itemize}[leftmargin=2em]
  \item Der \ISI-Kompositwert reicht von Frankreich~(0{,}236;
    Rang~27) bis Malta~(0{,}517; Rang~1).  Der EU-27-Mittelwert
    betr\"agt~0{,}344 bei einer Standardabweichung von~0{,}070.

  \item \begin{sloppypar}Die kovarianz\-basierte Varianz\-zerlegung
    ergibt: Die R\"ustungsachse tr\"agt~35{,}1\,\% zur
    Komposit\-varianz bei, die Logistik\-achse~34{,}3\,\%.
    Gemeinsam entfallen auf diese beiden
    Achsen~69{,}4\,\% der gesamten
    Querschnitts\-varianz.\end{sloppypar}

  \item Die \"ubrigen vier Achsen --- Finanzen, Energie, Technologie,
    kritische Rohstoffe --- tragen zusammen~30{,}6\,\% bei.  Ihre
    Werte sind im EU-Querschnitt vergleichsweise homogen verteilt.

  \item In 16 von 27~Mitgliedstaaten weist die R\"ustungsachse den
    h\"ochsten Einzelachsenwert auf; in den verbleibenden~11 ist es
    die Logistik\-achse.  Keine der \"ubrigen vier Achsen stellt in
    einem Mitgliedstaat den Spitzenwert.

  \item Das Entfernen der R\"ustungsachse aus dem Komposit verursacht
    die gr\"o\ss te Rangverschiebung
    (Spearman~$\rho = 0{,}833$ gegen\"uber dem Baseline-Ranking).
    Alle anderen Einzelachsen-Entfernungen ergeben Korrelationen
    \"uber~0{,}92.

  \item Unter Dirichlet-Gewichtsperturbation (10\,000~Ziehungen)
    betr\"agt die mittlere Rang\-ordnungs\-korrelation
    $\rho = 0{,}979$.  Die Rangordnung ist gegen\"uber moderaten
    Gewichts\-variationen stabil.

  \item \textbf{\"Osterreich} belegt Rang~8 (Kompositwert~0{,}383;
    moderat konzentriert).  H\"ochster Achsenwert:
    R\"ustung~(0{,}585), gefolgt von Logistik~(0{,}526).  Die
    Rangposition ist strukturell stabil --- unter keiner
    Einzelachsen-Perturbation verschiebt sich \"Osterreich um mehr
    als zwei R\"ange.
\end{itemize}


\newpage
% ─────────────────────────────────────────────────────────────────────────────
\section{Was der ISI misst}
\label{sec:was-der-isi-misst}
% ─────────────────────────────────────────────────────────────────────────────

Der \ISI ist ein zusammengesetzter Indikator, der die
\emph{Konzentration externer Lieferanten} erfasst.  F\"ur jeden
EU-Mitgliedstaat wird erhoben, inwieweit Importe in sechs strategisch
relevanten Bereichen von einer kleinen Anzahl externer Lieferanten
abh\"angen.  Der Index umfasst sechs gleichberechtigte Dimensionen:

\begin{table}[htbp]
\centering\small
\caption{Die sechs \ISI-Achsen und ihre Erfassungs\-bereiche.}
\label{tab:achsen}
\begin{tabular}{@{}lp{0.64\linewidth}@{}}
\toprule
\textbf{Achse} & \textbf{Erfassungsbereich} \\
\midrule
Finanzen &
  Grenz\"uber\-schreitende Bank- und
  Portfolio\-konzentration \\
Energie &
  Abh\"angigkeit von wenigen prim\"aren
  Energie\-lieferanten \\
Technologie &
  Import\-konzentration bei
  Hochtechnologie\-g\"utern \\
R\"ustung &
  Lieferanten\-konzentration bei R\"ustungs\-transfers und
  verteidigungs\-relevanten Importen \\
Kritische Inputs &
  Konzentration bei kritischen Rohstoffen und
  Vorleistungs\-g\"utern \\
Logistik &
  Abh\"angigkeit von wenigen Transport\-korridoren und
  Schifffahrts\-routen \\
\bottomrule
\end{tabular}
\end{table}

\begin{sloppypar}
Jeder Achsenwert basiert auf einem Herfindahl-Hirschman-\"ahnlichen
Konzentrations\-ma\ss, normiert auf das Einheits\-intervall
$[0,\,1]$.  Ein Wert von Null bedeutet vollst\"andige
Diversifikation; ein Wert nahe Eins bedeutet nahezu vollst\"andige
Abh\"angigkeit von einem einzigen externen Lieferanten.  Alle sechs
Achsen sind identisch skaliert.
\end{sloppypar}

\begin{sloppypar}
Der Kompositwert ist das \emph{gleichgewichtete arithmetische
Mittel} der sechs Achsenwerte.  Die Aggregation ist
deterministisch --- keine Sch\"atzung, keine latenten Faktoren,
keine Modell\-selektion.  Die Methodik ist vollst\"andig offen
dokumentiert: Paper~1~\cite{isi_paper1} legt den methodischen
Rahmen dar, Paper~2~\cite{isi_paper2} dokumentiert die empirische
Anwendung auf die EU-27.
\end{sloppypar}

\begin{sloppypar}
\textbf{Einordnung.}
Der \ISI misst keine ,,Souver\"anit\"at'' im normativen oder
verfassungs\-rechtlichen Sinne.  Er erfasst eine spezifische,
quantifizier\-bare strukturelle Eigenschaft: den Grad der
Konzentration externer Liefer\-beziehungen.  Hohe Konzentration
ist nicht automatisch negativ --- sie kann geographische N\"ahe,
Handels\-effizienz oder B\"undnis\-strukturen widerspiegeln.  Sie
impliziert jedoch, dass die Position des Landes gegen\"uber
St\"orungen in einer kleinen Anzahl von Liefer\-beziehungen
exponiert ist.
\end{sloppypar}


\newpage
% ─────────────────────────────────────────────────────────────────────────────
\section{Strukturelle Differenzierung innerhalb der EU-27}
\label{sec:differenzierung}
% ─────────────────────────────────────────────────────────────────────────────

\subsection{Varianzbeitr\"age der sechs Achsen}

\begin{sloppypar}
Der \ISI ist konzeptionell ein sechs\-dimensionaler Index.  Alle
sechs Achsen gehen gleichgewichtet in den Kompositwert ein und tragen
zum absoluten Konzentrations\-niveau jedes Landes bei.  Die
empirische Frage dieses Abschnitts ist eine andere: Welche Achsen
bestimmen die \emph{Unterschiede} zwischen den Mitgliedstaaten?
\end{sloppypar}

\begin{sloppypar}
Da alle sechs Achsen identisch auf das Intervall~$[0,\,1]$ normiert
sind und gleichgewichtet in den Komposit eingehen, ist die
kovarianz\-basierte Varianz\-zerlegung unmittelbar interpretierbar.
Die Querschnitts\-varianz der Kompositwerte wird \"uberwiegend durch
zwei Achsen bestimmt:
\end{sloppypar}

\begin{table}[htbp]
\centering\small
\caption{Varianzzerlegung des \ISI-Kompositwerts, EU-27.}
\label{tab:varianz}
\begin{tabular}{@{}lr@{}}
\toprule
\textbf{Achse} & \textbf{Anteil an Kompositvarianz} \\
\midrule
R\"ustung          & 35{,}1\,\% \\
Logistik           & 34{,}3\,\% \\
Kritische Inputs   & 20{,}9\,\% \\
Technologie        &  5{,}2\,\% \\
Energie            &  3{,}5\,\% \\
Finanzen           &  1{,}0\,\% \\
\midrule
R\"ustung + Logistik     & 69{,}4\,\% \\
\"Ubrige vier Achsen     & 30{,}6\,\% \\
\bottomrule
\end{tabular}
\end{table}

\begin{sloppypar}
Rund 64\,\% der Gesamtvarianz entfallen auf die Eigenvarianz
(Streuung jeder Achse \"uber die L\"ander), 36\,\% auf
Kreuz\-kovarianzen (gemeinsame Variation zwischen Achsen).
R\"ustung und Logistik dominieren beide Komponenten.
\end{sloppypar}

\textbf{Methodische Einordnung.}
Die Varianzzerlegung ist rein deskriptiv.  Sie beschreibt, welche
Achsen die Querschnittsvarianz statistisch bestimmen --- nicht,
welche Achsen politisch vorrangig sind.  Eine Achse kann politisch
von h\"ochster Relevanz sein und dennoch wenig zur
Querschnittsdifferenzierung beitragen, wenn alle Mitgliedstaaten
auf dieser Achse \"ahnliche Werte aufweisen.  Politische
Priorit\"at und statistische Differenzierung sind analytisch
getrennte Kategorien.

\begin{sloppypar}
Der Befund impliziert, dass strukturelle Heterogenit\"at innerhalb
der EU nicht gleichm\"a\ss ig \"uber alle strategischen Bereiche
verteilt ist, sondern sich auf wenige Dimensionen konzentriert.
Integrations\-politische Spannungen sind dort wahrscheinlicher
beobachtbar, wo nationale Profile stark auseinanderlaufen.
\end{sloppypar}

\subsection{Warum differenzieren Energie und Technologie weniger?}

\begin{sloppypar}
Energie\-abh\"angigkeit ist f\"ur viele Mitgliedstaaten von hoher
Relevanz.  Im EU-27-Querschnitt differenziert sie die L\"ander
jedoch weniger stark als R\"ustung und Logistik.  Der Grund ist
empirisch: Die Energie-Konzentrations\-werte streuen
vergleichsweise wenig.  Die meisten Mitgliedstaaten weisen
\"ahnliche Grade der Energie-Lieferanten\-konzentration auf.
Analoges gilt f\"ur die Technologie\-achse.  Beide Achsen tragen
zum absoluten Konzentrations\-niveau jedes Landes bei, erzeugen
aber im Vergleich zu R\"ustung und Logistik geringere Unterschiede
zwischen den L\"andern.
\end{sloppypar}

\subsection{Achsendominanz auf L\"anderebene}

\begin{sloppypar}
In 16~Mitgliedstaaten weist die R\"ustungsachse den h\"ochsten
Einzelachsenwert auf, in 11~die Logistik\-achse.  Keine der
\"ubrigen vier Achsen stellt in einem Mitgliedstaat den
Spitzenwert.  Dieses Muster beschreibt die empirische Verteilung
der Achsenmaxima --- es bedeutet nicht, dass der \ISI konzeptionell
auf zwei Achsen reduziert werden kann.  Der Index bleibt ein
sechs\-dimensionales Instrument; die beobachtete
Zweiachsen-Dominanz ist ein Befund \"uber die Daten, nicht \"uber
die Konstruktion.
\end{sloppypar}


% ─────────────────────────────────────────────────────────────────────────────
\section{\"Osterreich im EU-27-Kontext}
\label{sec:oesterreich}
% ─────────────────────────────────────────────────────────────────────────────

\"Osterreich belegt Rang~8 von~27 mit einem Kompositwert
von~0{,}383.  Es wird als ,,moderat konzentriert'' eingestuft ---
dieselbe Kategorie wie 23 der 27~Mitgliedstaaten.

\begin{table}[htbp]
\centering\small
\caption{\"Osterreichs Achsenprofil.}
\label{tab:oesterreich}
\begin{tabular}{@{}lrr@{}}
\toprule
\textbf{Achse} & \textbf{Wert} & \textbf{Achsenrang (v.\ 27)} \\
\midrule
R\"ustung          & 0{,}585 & 11. \\
Logistik           & 0{,}526 & 10. \\
Energie            & 0{,}456 &  3. \\
Kritische Inputs   & 0{,}319 &  7. \\
Technologie        & 0{,}261 &  5. \\
Finanzen           & 0{,}149 & 11. \\
\bottomrule
\end{tabular}
\end{table}

\begin{sloppypar}
\"Osterreichs h\"ochster Achsenwert entf\"allt auf
R\"ustung~(0{,}585), gefolgt von Logistik~(0{,}526).  Der Abstand
zwischen beiden betr\"agt~0{,}059 --- \"Osterreich weist kein
stark asymmetrisches Achsen\-profil auf.
\end{sloppypar}

\textbf{Energiekonzentration.}
\"Osterreichs drittst\"arkster Achsenwert ist Energie~(0{,}456),
Rang~3 in der EU.  Dies spiegelt die dokumentierte Abh\"angigkeit
von einer kleinen Zahl von Gaslieferanten wider.  Da die
Energie\-konzentration im EU-Querschnitt jedoch vergleichsweise
homogen verteilt ist, beeinflusst sie \"Osterreichs
Komposit-Ranking nur geringf"ugig.

\begin{sloppypar}
\textbf{Profilbalance.}
\"Osterreichs Variations\-koeffizient \"uber die sechs Achsen
betr\"agt~0{,}40 --- ein vergleichsweise ausgeglichenes
Achsen\-profil.  Zum Vergleich: Schweden (Rang~7;
Kompositwert~0{,}399) weist einen R\"ustungswert von~0{,}881 auf,
aber einen Finanzwert von nur~0{,}116 --- ein deutlich
asymmetrischeres Profil.
\end{sloppypar}

\textbf{Rangstabilit\"at.}
\begin{sloppypar}
Unter Leave-one-out-Analyse --- dem systematischen Entfernen jeder
einzelnen Achse --- verschiebt sich \"Osterreichs Rang um
h\"ochstens zwei Positionen.  Unter
Dirichlet-Gewichts\-perturbation (10\,000~Monte-Carlo-Ziehungen)
betr"agt "Osterreichs Rang-Standard\-abweichung ca.~1{,}1.  Das
95-\%-Konfidenz\-intervall umfasst die R\"ange~6 bis~10.
\"Osterreichs Position im oberen Mittelfeld ist strukturell robust.
\end{sloppypar}

\textbf{Nachbarschaft.}
\"Osterreich liegt zwischen Schweden (Rang~7; 0{,}399) und
Italien (Rang~9; 0{,}374).  Die Abst\"ande betragen~0{,}016 nach
oben und~0{,}009 nach unten.  Moderate Ver\"anderungen an einer
einzelnen Achse k\"onnten \"Osterreich um ein bis zwei
Rangpositionen verschieben; ein Sprung in die Top~5 oder unter
Rang~12 w\"urde substanzielle strukturelle Ver\"anderungen
erfordern.


\newpage
% ─────────────────────────────────────────────────────────────────────────────
\section{Robustheit der Ergebnisse}
\label{sec:robustheit}
% ─────────────────────────────────────────────────────────────────────────────

\begin{sloppypar}
Ein berechtigter Einwand gegen jeden zusammengesetzten Indikator
betrifft die Abh\"angigkeit der Ergebnisse von spezifischen
Methoden\-entscheidungen.  Drei Pr\"ufungen adressieren diesen
Einwand.
\end{sloppypar}

\subsection{Leave-one-out-Analyse}

\begin{sloppypar}
Wird jede Achse einzeln aus dem Komposit entfernt, ergeben sich
sechs alternative Rankings.  Die R\"ustungs-Entfernung erzeugt die
niedrigste Korrelation mit dem Baseline-Ranking
(Spearman~$\rho = 0{,}833$) --- dies best\"atigt, dass R\"ustung
die einzeln wirksamste Achse f\"ur die Rang\-differenzierung ist.
Logistik-Entfernung ergibt~$\rho = 0{,}926$.  Alle \"ubrigen
Entfernungen erzeugen Korrelationen \"uber~0{,}94.
\end{sloppypar}

Die Reihenfolge ist monoton: Achsen mit h\"oherem Varianzbeitrag
verursachen bei Entfernung gr\"o\ss ere Rangverschiebungen.

\subsection{Gewichtsstabilit\"at}

\begin{sloppypar}
Die Gleichgewichtung ist eine Entwurfs\-entscheidung.  Um deren
Konsequenzen zu testen, werden 10\,000~alternative
Gewichts\-vektoren aus einer symmetrischen Dirichlet-Verteilung
gezogen und der Kompositwert jeweils neu berechnet.  Die mittlere
Rang\-ordnungs\-korrelation \"uber alle Ziehungen
betr\"agt~0{,}979 --- die Rangordnung ist gegen\"uber moderaten
Gewichts\-variationen stabil.
\end{sloppypar}

\begin{sloppypar}
Die h\"ochste Rang\-volatilit\"at weisen die Niederlande auf
(Standard\-abweichung $\approx$~3{,}07~Rang\-positionen).  Diese
Volatilit\"at spiegelt das asymmetrische Achsen\-profil der
Niederlande wider: Der R\"ustungswert~(0{,}960) liegt weit \"uber
dem Niveau der \"ubrigen f\"unf Achsen.  St\"arkere Gewichtung der
R\"ustungsachse hebt den niederl\"andischen Rangplatz; schw\"achere
Gewichtung senkt ihn.  F\"ur die gro\ss e Mehrzahl der
Mitgliedstaaten ist das Ranking auf ein bis zwei Positionen stabil.
\end{sloppypar}

\subsection{Korrelationsstruktur}

\begin{sloppypar}
Die sechs Achsen sind weitgehend unabh\"angig voneinander.  Von
15~paarweisen Korrelationen sind nach Bonferroni-Korrektur nur
zwei statistisch signifikant: Energie--Technologie und Kritische
Inputs--Logistik.  Die Finanzachse ist zu allen anderen Achsen
nahezu orthogonal (h\"ochster $|t|$-Wert: 1{,}11).
\end{sloppypar}

Diese Unabh\"angigkeit ist analytisch vorteilhaft: Die
Varianzzerlegung spiegelt genuinen strukturellen Gehalt wider,
nicht statistische Redundanz.  Die geringe
Multikollinearit\"at verhindert eine k\"unstliche
Varianzaufbl\"ahung einzelner Achsen.


% ─────────────────────────────────────────────────────────────────────────────
\section{Einordnung und analytische Implikationen}
\label{sec:implikationen}
% ─────────────────────────────────────────────────────────────────────────────

Der \ISI schreibt keine Politik vor.  Er liefert eine empirische
Messgrundlage.

\begin{sloppypar}
\textbf{Querschnitts\-varianz und absolute Konzentration.}
R\"ustung und Logistik bestimmen die Querschnitts\-varianz.  Daraus
folgt \emph{nicht}, dass die \"ubrigen Achsen f\"ur einzelne
Mitgliedstaaten irrelevant sind.  Ein Land kann auf der
Energieachse einen hohen absoluten Konzentrations\-wert aufweisen,
ohne dass dieser Wert die L\"ander im Querschnitt voneinander
unterscheidet.  Politische Priorit\"at und statistische
Differenzierung sind analytisch getrennte Kategorien.
\end{sloppypar}

\textbf{Empirische Fundierung strategischer Debatten.}
\begin{sloppypar}
Der Begriff ,,strategische Autonomie'' ist in der
EU-Politikdebatte allgegenw\"artig, aber selten an vergleichbare,
quantitative Daten gekn\"upft.  Der \ISI bietet einen solchen
Ankerpunkt.  Die Varianz\-zerlegung zeigt, dass
R\"ustungs\-beschaffung und Logistik\-korridore die gr\"o\ss ten
Unterschiede zwischen den Mitgliedstaaten erzeugen --- Bereiche,
die in der \"offentlichen Debatte gegen\"uber Energie und
Technologie vergleichsweise wenig Aufmerksamkeit erhalten.
\end{sloppypar}

\begin{sloppypar}
\textbf{Energie und Technologie.}
Energie\-abh\"angigkeit ist f\"ur zahlreiche Mitgliedstaaten von
hoher Bedeutung.  Im EU-27-Querschnitt differenziert sie die
L\"ander jedoch weniger stark als R\"ustung und Logistik, weil die
Energie-Konzentrations\-werte vergleichsweise gleichm"a\ss ig
verteilt sind.  Analoges gilt f\"ur Technologie.  Beide Achsen
erfordern eigenst\"andige Analyse; ihr geringer
Streuungs\-beitrag schm\"alert ihre sachliche Bedeutung nicht.
\end{sloppypar}

\begin{sloppypar}
\textbf{\"Osterreich.}
Rang~8 platziert \"Osterreich im oberen Mittelfeld der externen
Konzentration.  Das vergleichsweise ausgeglichene Achsen\-profil
weist weder eine singul\"are Konzentrations\-spitze noch eine
au\ss ergew"ohnliche Diversifikations\-st"arke auf.
\end{sloppypar}

\begin{sloppypar}
\textbf{Messung, nicht Bewertung.}
Ein hoher Kompositwert bedeutet nicht, dass ein Land ,,schlecht
abschneidet''.  Maltas Rang~1 spiegelt Gr\"o\ss e und
geographische Lage wider, nicht politisches Versagen.
Frankreichs Rang~27 spiegelt die strukturelle Diversifikation
einer gro\ss en, gut vernetzten Volkswirtschaft wider.  Der Index
misst Konzentration --- was die Politik daraus ableitet, ist eine
separate Frage.
\end{sloppypar}


% ─────────────────────────────────────────────────────────────────────────────
\section{Limitationen}
\label{sec:limitationen}
% ─────────────────────────────────────────────────────────────────────────────

\textbf{Re-Export-Verzerrung.}
\begin{sloppypar}
Handelsstatistiken erfassen bilaterale Str\"ome, identifizieren
aber nicht immer den tats\"achlichen Ursprung der G\"uter.
L\"ander, die als Handels\-drehscheiben fungieren (etwa die
Niederlande oder Belgien), erscheinen m\"oglicherweise
diversifizierter als sie tats\"achlich sind.
\end{sloppypar}

\begin{sloppypar}
\textbf{Verst\"arkungseffekt bei kleinen Volkswirtschaften.}
Kleine, offene Volkswirtschaften mit begrenzter Zahl an
Handelspartnern erzeugen mechanisch h\"ohere
Konzentrations\-werte.  Malta und Zypern belegen Rang~1 und~2
teilweise, weil Geographie und Gr\"o\ss e ihre Bezugsoptionen
einschr\"anken.  Dies ist ein reales Merkmal ihrer Situation, kein
Messfehler --- es sollte jedoch entsprechend eingeordnet werden.
\end{sloppypar}

\begin{sloppypar}
\textbf{Querschnitts\-limitation.}
Der Datensatz bildet den Jahrgang~2024 ab.
Lieferanten\-konzentration kann sich durch politische
Entscheidungen, Markt\-entwicklungen oder geopolitische Ereignisse
ver\"andern.  Ein einzelner Querschnitt kann keine Dynamiken
erfassen.  K\"unftige \ISI-Jahrg\"ange werden Trend\-analysen
erm\"oglichen.
\end{sloppypar}

\textbf{Gleichgewichtung.}
\begin{sloppypar}
Die Aggregation behandelt alle sechs Achsen als gleich wichtig.
Das ist eine normative Setzung.  Ein Entscheidungs\-tr\"ager, der
Energie\-abh\"angigkeit als vorrangig erachtet, w\"urde andere
Gewichte setzen und ein anderes Ranking erhalten.  Die
Dirichlet-Perturbations\-analyse zeigt, dass moderate Abweichungen
von der Gleichgewichtung die Ergebnisse nicht grunds\"atzlich
ver\"andern.
\end{sloppypar}

\vfill
\begin{center}
\rule{0.4\linewidth}{0.4pt}
\end{center}
\smallskip
\noindent\textit{Dieses Briefing basiert auf der ISI Paper Series
(Papers~1--2), erstellt von der Research Unit des International
Sovereignty Institute.  Die zugrunde liegenden Daten und die
vollst\"andige Methodik sind \"offentlich dokumentiert.}
